\chapter{The Log-Structured Alternative}
\label{ch:lsm}

A B+ tree updates pages \emph{in place}: an insert finds the right leaf and
changes it. The other great family of storage engines never updates in place.
A \term{log-structured merge tree} (LSM tree), introduced by O'Neil et al.\
\cite{lsm}, turns every write into a sequential append and pays for it later,
in the background. RocksDB, LevelDB, Cassandra, and many modern stores are LSM
trees. Knowing why OntoDB does \emph{not} use one teaches as much as knowing why
it uses a B+ tree.

\section{How an LSM tree works}

\begin{enumerate}
\item \textbf{Memtable.} Writes go to an in-memory sorted structure (often a skip list), and to a write-ahead log for durability.
\item \textbf{Flush.} When the memtable is full, it is written to disk as an immutable sorted file, an \term{SSTable}, in one sequential write.
\item \textbf{Compaction.} SSTables accumulate. A background process merges them into larger SSTables, dropping overwritten and deleted entries. In \term{leveled} compaction, level $L+1$ is about ten times larger than level $L$.
\item \textbf{Reads} check the memtable, then SSTables from newest to oldest, and stop at the first hit.
\item \textbf{Deletes} write a \term{tombstone}, a marker that hides older versions until compaction removes both.
\end{enumerate}

A read that misses must look in many files. \term{Bloom filters}, one per
SSTable, answer ``is this key possibly here?'' in memory with no false
negatives, so most files are skipped without a read.

\section{Three amplifications}

Every storage engine trades three costs against each other.

\begin{description}[leftmargin=1.2em]
\item[Write amplification] Bytes written to the device per byte of user data. An LSM tree rewrites data at every compaction; a B+ tree rewrites a whole page to change a few bytes.
\item[Read amplification] Pages read per lookup. A B+ tree reads one path; an LSM tree may consult several levels.
\item[Space amplification] Bytes stored per byte of live data. LSM trees hold obsolete versions until compaction; B+ trees hold half-empty pages.
\end{description}

\begin{table}[ht]
\centering
\begin{tabular}{@{}lll@{}}
\toprule
& B+ tree & LSM tree \\
\midrule
Writes & random, in place & sequential, batched \\
Point reads & one root-to-leaf path & memtable, then levels (Bloom filters help) \\
Range scans & walk linked leaves & merge iterators across levels \\
Deletes & remove now, merge pages & tombstone now, remove at compaction \\
Background work & none required & compaction, always \\
Best for & read-heavy, range-heavy & write-heavy, ingest-heavy \\
\bottomrule
\end{tabular}
\caption{The two families, in broad strokes.}
\end{table}

\begin{keyidea}
An LSM tree buys fast writes with deferred work. A B+ tree pays for each write
immediately, and gets predictable reads and no background compaction in return.
\end{keyidea}

\begin{ontodb}[: why a B+ tree]
OntoDB's workload is a SPARQL engine's: a bulk load, then queries that are joins
of prefix range scans. Range scans over sorted permutations are a B+ tree's
best case, and bulk loading builds full, sequential leaves without any
compaction. A delete in OntoDB's tree shrinks the tree in place, inside
\code{BPlusTree::Remove}; there is no tombstone for a reader to skip and no
background job that rewrites sorted runs.

The LSM ideas still appear in OntoDB, in smaller roles. Its persistent
databases (\cref{ch:durability}) are log-plus-snapshot: every change is appended
to a log, and a checkpoint rewrites the whole state as one sorted image, which
is a compaction with a single level. And Bloom filters are a natural fit for
``has this concept been modified recently?'' checks over a log.
\end{ontodb}

\begin{checkyourself}
\item A workload inserts 10 million triples per hour and runs few queries. Which family fits better, and what does the other one pay?
\item Why does an LSM delete write a tombstone instead of removing the key?
\item A Bloom filter says ``maybe present'' for a key that is absent. What does that cost an LSM read? Can it say ``absent'' for a present key?
\item In what sense is OntoDB's snapshot-plus-log a one-level LSM tree?
\end{checkyourself}

\section*{Further reading}
O'Neil et al.\ \cite{lsm}; Petrov \cite{petrov}, chapter 7; Bloom
\cite{bloom}.
